\documentclass[10pt,a4paper]{amsart}
\usepackage[margin=19mm]{geometry}
\usepackage{amsmath,amssymb,mathtools}
\usepackage[scr=rsfs,cal=euler]{mathalfa}
\usepackage{xfrac}
\usepackage[unicode,hidelinks,bookmarks=false]{hyperref}
\newcommand{\ie}{i.e.\ }
\newcommand{\cf}{cf.\ }
\newcommand{\resp}{respectively\ }
\newcommand{\Subsection}{\subsection}
\providecommand{\nobreakdash}{\nobreak\hbox{-}}
\setlength{\emergencystretch}{2em}
\multlinegap0pt
\usepackage{amssymb}
\usepackage{enumerate}
\usepackage{mathtools}
\usepackage{tikz}
\newtheorem{thm}{Theorem}
\title{Thurston norm and the Euler class}
\author{Mehdi Yazdi}
\date{Source: arXiv:2604.08096v1 (CC BY 4.0)}
\begin{document}
\maketitle

\noindent\emph{Source and licence.} Thurston norm and the Euler class. Version arXiv:2604.08096v1 (CC BY 4.0), distributed by arXiv under the Creative Commons Attribution 4.0 International licence, http://creativecommons.org/licenses/by/4.0/, as recorded in the arXiv licence metadata for this exact version and shown on the abstract page https://arxiv.org/abs/2604.08096v1. Full licence notice: \texttt{licenses/arxiv-2604.08096-CC-BY-4.0.txt}.\par\smallskip

\noindent\textit{Editorial scope.} Counted unit: the article's thm vertices (source0.tex lines 115--118). The article's complete original proof of it is delivered with it (source0.tex lines 610--636, one \texttt{proof} environment of 27 lines, which the article places away from the statement). No mathematics is written, repaired or re-derived: the statement and the proof are the article's own lines, in the article's own notation, and no retained line is substituted or re-flowed.  Retained uncounted as the source-local context the proof needs: lines 128--136; lines 601--606.  Nothing was omitted from any retained span, so the package carries no omission record.  The retained lines cite 1 works, whose entries are transcribed verbatim from the article's own bibliography source in the e-print and delivered with short editorial labels recorded in the item's reference mapping.  No retained line contains a figure environment, an image-inclusion command or drawing code, so no image is delivered and the package declares no asset dependency.  Editorial formatting only, in addition: an AMS \texttt{amsart} preamble that restates the article's own declarations and macros used by the retained lines, namely the environments \texttt{thm} on separate counters, together with the standard packages those lines need.  The retained environments print in this document as Theorem 1 in order of appearance, whereas the article prints the same items with the numbering of its own full document.  The complete native source of the article as distributed in the arXiv e-print for this exact version is bundled unchanged as \texttt{sources/198135/source0.tex}.
	 \begin{thm}[Fully Marked Surface Theorem]
	 	Let $M$ be a closed hyperbolic 3-manifold, $\mathcal{F}$ be a taut foliation of $M$, and $S$ be an algebraically fully marked surface in $M$. There is a taut foliation $\mathcal{G}$ of $M$ and an embedded surface $S'$ in $M$ such that 
	 	\begin{enumerate}
	 		\item $S'$ is homologous to $S$; 
	 		\item the oriented tangent plane fields of $\mathcal{F}$ and $\mathcal{G}$ are homotopic through plane fields; and 
	 		\item $S'$ is a  union of compact leaves of $\mathcal{G}$.
	 	\end{enumerate} 
	 	\label{thm:fully marked surface}
	 \end{thm}

	\begin{thm}[Liu]
		Let $M$ be a closed oriented hyperbolic 3-manifold. Denote by $B_x(M)$ the unit ball of the Thurston norm of $M$, and by $B_{x^*}(M)$ the unit ball of the dual norm. Let $F \subset \partial B_{x^*}(M)$ and $F^\wedge \subset \partial B_x(M)$ be a dual pair of closed faces.  Suppose that $w \in H^2(M ; \mathbb{R})$ is a rational point in the interior of $F$, and $\Sigma \in H_2(M ; \mathbb{Z})$ is a primitive homology class such that $\Sigma$ lies in the interior of $F^\wedge$. 
		There exists some finite cyclic covering $p \colon \tilde{M} \rightarrow M$ dual to $\Sigma$, and some transversely oriented taut foliation $\tilde{\mathcal{F}}$ on $\tilde{M}$ such that the following equality holds in $H^2(M ; \mathbb{R})$
		\[ p!(e(\mathcal{F})) = [\tilde{M}: M] \cdot w. \]
		\label{thm:Liu-virtual-taut foliation-projection}
	\end{thm}

	\begin{thm}[Euler Class One for Vertices]\index{Euler class!Euler class one theorem for vertices}
		Let $M$ be a compact orientable irreducible 3-manifold with boundary a (possibly empty) union of tori, and with positive first Betti number. Every vertex of the unit ball of the dual Thurston norm is realised as the Euler class of some taut foliation on $M$. 
		\label{thm:Gabai-euler class one for vertices}
	\end{thm}

\begin{proof}[Sketch of the proof of Theorem \ref{thm:Liu-virtual-taut foliation-projection}]
	
	Let $v_1, \cdots, v_n \in F$ be the set of vertices of $F$. By the Euler class one theorem for vertices (Theorem \ref{thm:Gabai-euler class one for vertices}), for each $v_i$, there exists a taut foliation $\mathcal{F}_i$ on $M$ with Euler class $e(\mathcal{F}_i) = v_i$. Since $\Sigma$ is in the cone over $F^\wedge$, we have the following equality for every $i$
	\begin{eqnarray}
		\langle e(\mathcal{F}_i) , \Sigma \rangle = x(\Sigma), 
		\label{eq:fully marked}
	\end{eqnarray}
	where $x$ denotes the Thurston norm.
	By the fully marked surface theorem (Theorem \ref{thm:fully marked surface}), possibly after replacing $\mathcal{F}_i$ with new foliations but preserving Equality (\ref{eq:fully marked}), we may assume that for each $i$ there exists a closed (possibly disconnected) embedded oriented surface $S_i \subset M$ such that $[S_i] = \Sigma$, and $S_i$ is a union of compact leaves of $\mathcal{F}_i$. In general $S_i$ would be in different isotopy classes and possibly intersect each other, but to simplify the exposition and to bring the main ideas across here we make the assumption that $\Sigma$ is the unique norm-minimising surface in its homology class. For the complete proof without this assumption see Liu \cite{liu2024criterion}. Because of our simplifying assumption, $S_i = \Sigma$ up to isotopy, and so $\Sigma$ is a common leaf of $\mathcal{F}_1, \cdots, \mathcal{F}_n$. 
	
	Since $w \in \mathrm{int} (F) $ we can write it as 
	\[ w =  \frac{a_1 v_1 + \cdots + a_n v_n}{a_1 + \cdots + a_n},\]
	for positive integers $a_1 , \cdots, a_n$. Consider the foliation $\mathcal{F}_i \setminus \setminus \Sigma$ obtained by cutting $\mathcal{F}_i$ along the compact leaf $\Sigma$. There are two copies of $\Sigma$ in the boundary of $\mathcal{F}_i \setminus \setminus \Sigma$, denote them by $\Sigma_+$ and $\Sigma_-$. Consider a cyclic cover $\tilde{M}$ of $M$ dual to $\Sigma$ and of degree $(a_1 + \cdots + a_n)$, obtained by stacking $a_i$ copies of $\mathcal{F}_i \setminus \setminus \Sigma$ together for all $1 \leq i \leq n$, glued along the copies of $\Sigma$ such that $\Sigma_+$ in one copy is glued to $\Sigma_-$ in the adjacent copy. By construction $\tilde{M}$ comes equipped with a foliation $\mathcal{F}$ obtained by gluing a suitable number of copies of the foliations $\mathcal{F}_i \setminus \setminus \Sigma$ together. It is shown in \cite[Lemma 4.3]{liu2024criterion} that the Euler class $e(\mathcal{F})$ satisfies 
	\begin{eqnarray}
		p!(e(\mathcal{F})) = a_1 v_1 + \cdots + a_n v_n,
		\label{eq:euler class of medley construction}
	\end{eqnarray}
	or equivalently
	\[  p!(e(\mathcal{F}))  = (a_1 + \cdots + a_n) w = [\tilde{M} : M] \cdot w. \]
	We now justify Equality (\ref{eq:euler class of medley construction}) by a topological argument. Both sides of Equality (\ref{eq:euler class of medley construction}) are elements of $H^2(M ; \mathbb{R})$, and so it is enough to show that for every embedded oriented surface $S$ in $M$, the evaluations of both sides on the homology class $[S]$ are equal to each other. By definition of the transfer map we have 
	\[ \langle p!(e(\mathcal{F}))  , [S] \rangle = \langle e(\mathcal{F}) , [p^{-1}(S)] \rangle . \]
	So it is enough to show that 
	\begin{align*} \langle e(\mathcal{F}) , [p^{-1}(S)] \rangle  &= \langle  a_1 v_1 + \cdots + a_n v_n , [S] \rangle =  a_1 \langle e(\mathcal{F}_1) , [S] \rangle + \cdots + a_n  \langle e(\mathcal{F}_n) , [S] \rangle.
	\end{align*}
	Isotope $S$ so that it is in general position with respect to $\Sigma$ and intersects $\Sigma$ in a union of simple closed curves. Then $\langle e(\mathcal{F}_i) , [S] \rangle $ is the obstruction for finding a section of $T \mathcal{F}_i$ over $S$. Define a common section of $T \mathcal{F}_i$ over $\Sigma \cap S$ as any non-vanishing vector field $s$ in $T \Sigma \cap T S$. Let $M'$ be the manifold obtained by cutting $M$ along $\Sigma$, and $\mathcal{F}_i'$ be the foliation on $M'$ obtained by cutting $\mathcal{F}_i$ along the compact leaf $\Sigma$.  Therefore $\langle e(\mathcal{F}_i) , [S] \rangle $ is equal to the obstruction for extending the section $s$ to a section of $T \mathcal{F}_i'$ over $S'$. Now the foliation $\mathcal{F}$ is obtained by stacking together $a_i$ copies of $\mathcal{F}_i'$, glued along copies of $\Sigma$. The section $s$ lifts to a section $\tilde{s}$ of $\mathcal{F}$ on $p^{-1}(\Sigma \cap S) = p^{-1}(\Sigma) \cap p^{-1}(S) $, and $\langle e(\mathcal{F}) , [p^{-1}(S)] \rangle $ is the obstruction for extending $\tilde{s}$ to a section over $p^{-1}(S)$. Now the latter obstruction is the sum of obstructions coming from extending the section $\tilde{s}$ to the part of $p^{-1}(S)$ lying inside $a_i$ copies of $\mathcal{F}_i'$ for $1 \leq i \leq n$, which by the previous discussion is equal to $ a_i \langle e(\mathcal{F}_i) , [S] \rangle $. The equality follows. 

\end{proof}

\begin{thebibliography}{99}
\bibitem[liu202]{liu2024criterion} Yi~Liu. \newblock A criterion for virtual {E}uler class one. \newblock {\em arXiv preprint arXiv:2411.11492}, 2024.
\end{thebibliography}
\end{document}
